\scanpage{39}
\margin{38}
que toute fonction rationnelle definie en $x$ soit definie en $y$. \deleted{Alors $\underline O_y$ domine} La condition est de toutes façons necessaire, prouvons qu'elle est suffisante (utilisant la separation). Supposons $\underline O_x\quad\underline O_y$, donc (\deleted{\unsure{l}}lemme 1) il existe un ideal premier $\underline p$de $\underline O_x$ tel que $\underline O_y$ domine $\underline O_{x\,\underline p}$. Or il existe un $z\quad X$ tel que $x\quad\bar z$, $\underline O_z=\underline O_{x\,\underline p}$ (\unsure{reference}
donc $\underline O_z$ et $\underline O_y$ sont apparentes, donc $z=y$ cqfd.

On peut reenoncer ce corollaire ainsi: si $x\notin\bar y$, il existe une fonction rationnelle $f$ sur $X$ telle que $f$ soit definie en $x$ et non en $y$ i.e. telle que $D(f)$ contienne $x$ et ne rencontre pas $\bar y$. Si maintenant $X$ est un espace noetherien, toute partie fermee de $X$ est reunion finie de fermes irreductibles, i.e. de la forme $\bar y$, d'ou:

\uline{Corollaire} 3 Soit $X$ un schema integre noetherien (il suffit de dire: qui est un espace topologique noeth.), alors les $D(f)$ forment une base de la topologie de $X$.

On voit ainsi comment la topologie de $X$, identifie a un ensemble de sous-anneaux locaux de $R$, peut se recuperer en termes de cette derniere structure; de façon equivalente, on peut aussi remarquer que la topologie de $X$ est definie entermes de sastructure d'ordre, un \unsure{ferme} de $X$ etant forme des elements de $X$ majores par un au moins parmi les elements d'une partie finie de $X$ (la majoration etant definie \deleted{\unsure{sont}} ici en termes d'inclusions \deleted{de parties \unsure{de $R$}} d'anneaux locaux). D'autre part, le faisceau $\underline O_X$ lui aussi est recupere, puisque
\[
\mathrm I(U,\underline O_X)=\underline O_x\quad(x\quad U)
\]
Plus bas, nous donnerons quelques indications plus precis\deleted{\unsure{s}}es.

\uline{Proposition} 3 Soit $f:X\quad Y$ un morphisme dominant de schemas integres, \deleted{\unsure{de corps}} soit $L$ ($K$) le corps de $X$ ($Y$). Identifions $K$ a un sous-corps de $L$, et les anneaux locaux de $X$ ($Y$) a des sous-anneaux de $L$ ($K$). Alors pour tout $x\quad X$, $\underline O_{f(x)}$ est l'unique anneau local de $Y$ domine par $\underline O_x$.

Il est en effet domine, par definition d'un morphisme, et il est unique puisque deux anneaux locaux de $Y$ apparentes sont identiques.

\uline{Proposition} 4 Soit $f$ un morphisme \uline{separe} d'immersion locale (resp. d'isomorphisme local) de $X$ dans $Y$, avec \deleted{$Y$} $X$ \uline{irreductible}. Alors $f$ est une immersion (resp. une immersion ouverte).

Il suffit de montrer que $f$ est un homeomorphisme dans (ce qui impliquera facilement, compte tenu de l'hypothese, que $f(X)$ est localement ferme dans $Y$). On peut evidemment supposer $X$ et $Y$ \uline{reduits}, et on voit qu'on peut remplacer $Y$ par l'adherence de $f(X)$, donc on peut supposer $f$ dominant. Alors \deleted{$Y$ sera aussi} $Y$ sera aussi irreductible, donc $X$ et $Y$ sont integres. Avec les notations de prop.3, on aura pour tout $x\quad X$
